\renewcommand{\Anno}{1997}

\renewcommand{\NumeroProblema}{3}

\renewcommand{\Testo}{
La seguente figura rappresenta una scacchiera con le caselle 
	colorate nel modo usuale.
\begin{center}
		\psset{unit=5mm}
		\psset{linewidth=0.6pt}
		\newrgbcolor{scacchiera}{0.4 1 0.6}
	
		\pspicture(0,0)(8,8)
	
		\multips(0,0)(2,0){4}{\pspolygon*[linecolor=scacchiera](0,0)(1,0)(1,1)(0,1)(0,0)}
		\multips(1,1)(2,0){4}{\pspolygon*[linecolor=scacchiera](0,0)(1,0)(1,1)(0,1)(0,0)}
		\multips(0,2)(2,0){4}{\pspolygon*[linecolor=scacchiera](0,0)(1,0)(1,1)(0,1)(0,0)}
		\multips(1,3)(2,0){4}{\pspolygon*[linecolor=scacchiera](0,0)(1,0)(1,1)(0,1)(0,0)}
		\multips(0,4)(2,0){4}{\pspolygon*[linecolor=scacchiera](0,0)(1,0)(1,1)(0,1)(0,0)}
		\multips(1,5)(2,0){4}{\pspolygon*[linecolor=scacchiera](0,0)(1,0)(1,1)(0,1)(0,0)}
		\multips(0,6)(2,0){4}{\pspolygon*[linecolor=scacchiera](0,0)(1,0)(1,1)(0,1)(0,0)}
		\multips(1,7)(2,0){4}{\pspolygon*[linecolor=scacchiera](0,0)(1,0)(1,1)(0,1)(0,0)}
	
		\multido{\n=0+1}{9}{
		\psline{-}(\n,0)(\n,8)
		\psline{-}(0,\n)(8,\n)
		}
	
		\endpspicture
\end{center}

	Dalla posizione iniziale si possono ottenere altre configurazioni
	della scacchiera scambiando tra loro due righe o due colonne
	qualsiasi, ripetendo l'operazione quante volte si vuole.

	Determinare quali delle seguenti configurazioni 
\begin{center}
		\hfill
% 		\psset{unit=5mm}
% 		\psset{linewidth=0.6pt}
% 		\newrgbcolor{scacchiera}{0.4 1 0.6}
% 		\pspicture(0,0)(8,8)
% 	
% 		\multips(1,0)(2,0){4}{\pspolygon*[linecolor=scacchiera](0,0)(1,0)(1,1)(0,1)(0,0)}
% 		\multips(0,1)(2,0){4}{\pspolygon*[linecolor=scacchiera](0,0)(1,0)(1,1)(0,1)(0,0)}
% 		\multips(1,2)(2,0){4}{\pspolygon*[linecolor=scacchiera](0,0)(1,0)(1,1)(0,1)(0,0)}
% 		\multips(0,3)(2,0){4}{\pspolygon*[linecolor=scacchiera](0,0)(1,0)(1,1)(0,1)(0,0)}
% 		\multips(1,4)(2,0){4}{\pspolygon*[linecolor=scacchiera](0,0)(1,0)(1,1)(0,1)(0,0)}
% 		\multips(0,5)(2,0){4}{\pspolygon*[linecolor=scacchiera](0,0)(1,0)(1,1)(0,1)(0,0)}
% 		\multips(1,6)(2,0){3}{\pspolygon*[linecolor=scacchiera](0,0)(1,0)(1,1)(0,1)(0,0)}
% 		\multips(0,7)(2,0){3}{\pspolygon*[linecolor=scacchiera](0,0)(1,0)(1,1)(0,1)(0,0)}
% 	
% 		\multido{\n=0+1}{9}{
% 		\psline{-}(\n,0)(\n,8)
% 		\psline{-}(0,\n)(8,\n)
% 		}
% 	
% 		\endpspicture
		\psset{unit=5mm}
		\psset{linewidth=0.6pt}
		\newrgbcolor{scacchiera}{0.4 1 0.6}
		\pspicture(0,0)(8,8)
	
		\pspolygon*[linecolor=scacchiera](0,8)(0,4)(8,4)(8,0)(5,0)(5,8)(0,8)
	
		\multido{\n=0+1}{9}{
		\psline{-}(\n,0)(\n,8)
		\psline{-}(0,\n)(8,\n)
		}
	
		\endpspicture
		\hfill
		\pspicture(0,0)(8,8)
	
		\pspolygon*[linecolor=scacchiera](0,8)(0,4)(1,4)(1,3)(2,3)(2,2)(3,2)(3,1)
		(4,1)(4,0)(8,0)(8,1)(7,1)(7,2)(6,2)(6,3)(5,3)(5,4)(4,4)(4,5)(3,5)
		(3,6)(2,6)(2,7)(1,7)(1,8)(0,8)
		\pspolygon*[linecolor=scacchiera](5,8)(5,7)(6,7)(6,6)(7,6)(7,5)(8,5)(8,8)(6,8)
	
		\multido{\n=0+1}{9}{
		\psline{-}(\n,0)(\n,8)
		\psline{-}(0,\n)(8,\n)
		}
	
		\endpspicture
		\hfill
		\pspicture(0,0)(8,8)
	
		\pspolygon*[linecolor=scacchiera](0,0)(2,0)(2,2)(0,2)(0,0)
		\pspolygon*[linecolor=scacchiera](0,6)(2,6)(2,8)(0,8)(0,6)
		\pspolygon*[linecolor=scacchiera](6,0)(8,0)(8,2)(6,2)(6,0)
		\pspolygon*[linecolor=scacchiera](6,6)(8,6)(8,8)(6,8)(6,6)
		\pspolygon*[linecolor=scacchiera](2,2)(6,2)(6,6)(2,6)(2,2)
	
		\multido{\n=0+1}{9}{
		\psline{-}(\n,0)(\n,8)
		\psline{-}(0,\n)(8,\n)
		}
	
		\endpspicture
		\hfill\mbox{}
\end{center}
	si possono ottenere e quali non si possono ottenere.
}

\renewcommand{\Risposta}{ 
La scacchiera di sinistra e quella centrale non sono ottenibili,
quella di destra \`{e} ottenibile.
}

\renewcommand{\Aiutino}{
Procedere per tentativi per familiarizzare con il gioco e scoprire 
invarianti.
\begin{itemize}
	\item  Per la scacchiera di sinistra: quante caselle colorate ci 
	devono essere in ogni riga ed ogni colonna dopo ogni mossa?

	\item  Per la scacchiera centrale: quanti tipi diversi di riga o di 
	colonna possono/devono essere presenti contemporaneamente sulla scacchiera 
	al termine di ogni mossa?

	\item  Provare, provare, provare.
\end{itemize}
}

\renewcommand{\SoluzioneUfficiale}{
Consideriamo i seguenti due invarianti del problema.
\begin{itemize}
	\item  Primo invariante: dopo qualunque successione di mosse, in ogni 
	riga ci sono sempre 4 caselle bianche e 4 caselle colorate (lo stesso 
	vale ovviamente per le colonne). Se infatti si fa uno scambio di 
	righe, queste rimangono le stesse del passo precedente, mentre 
	quando si fa uno scambio di colonne, quella che pu\`{o} cambiare 
	\`{e} la disposizione delle caselle colorate all'interno di ogni 
	riga, ma non il loro numero.
	
	Questo invariante basta per escludere che si possa ottenere la 
	scacchiera di sinistra, nella quale sono presenti righe con 3 o 5 
	caselle colorate.

	\item  Secondo invariante: dopo qualunque successione di mosse, 
	sulla scacchiera le righe sono uguali a 4 a 4 (stesso discorso per 
	le colonne). Se infatti si fa uno scambio di righe, queste rimangono 
	le stesse del passo precedente, se invece si fa uno scambio di 
	colonne, le 4 righe che erano uguali prima dello spostamento 
	cambieranno allo stesso modo, rimanendo dunque uguali al termine 
	della mossa.
	
	Questo basta per escludere che si possa ottenere la scacchiera 
	centrale, nella quale le righe sono addirittura tutte diverse.
\end{itemize}

La scacchiera di destra si pu\`{o} invece ottenere, ad esempio con il
procedimento descritto nella seguente figura (vengono eseguiti 2
scambi di colonne e 2 scambi di righe).
\begin{center}
	\hfill
	\psset{unit=5mm}
	\newrgbcolor{scacchiera}{0.4 1 0.6}
	\pspicture(0,-0.8)(8,8)
	
	\multips(0,0)(2,0){4}{\pspolygon*[linecolor=scacchiera](0,0)(1,0)(1,1)(0,1)(0,0)}
	\multips(1,1)(2,0){4}{\pspolygon*[linecolor=scacchiera](0,0)(1,0)(1,1)(0,1)(0,0)}
	\multips(0,2)(2,0){4}{\pspolygon*[linecolor=scacchiera](0,0)(1,0)(1,1)(0,1)(0,0)}
	\multips(1,3)(2,0){4}{\pspolygon*[linecolor=scacchiera](0,0)(1,0)(1,1)(0,1)(0,0)}
	\multips(0,4)(2,0){4}{\pspolygon*[linecolor=scacchiera](0,0)(1,0)(1,1)(0,1)(0,0)}
	\multips(1,5)(2,0){4}{\pspolygon*[linecolor=scacchiera](0,0)(1,0)(1,1)(0,1)(0,0)}
	\multips(0,6)(2,0){4}{\pspolygon*[linecolor=scacchiera](0,0)(1,0)(1,1)(0,1)(0,0)}
	\multips(1,7)(2,0){4}{\pspolygon*[linecolor=scacchiera](0,0)(1,0)(1,1)(0,1)(0,0)}
	
	\multido{\n=0+1}{9}{
	\psline{-}(\n,0)(\n,8)
	\psline{-}(0,\n)(8,\n)
	\psarc[linewidth=0.6\pslinewidth]{<->}(2,0){0.7}{200}{340}
	\psarc[linewidth=0.6\pslinewidth]{<->}(6,2){2.7}{235}{305}
	}
	
	\endpspicture
	\hfill
	\pspicture(0,-0.8)(9,8)
	
	\multips(0,0)(0,2){4}{\pspolygon*[linecolor=scacchiera](0,0)(2,0)(2,1)(0,1)(0,0)}
	\multips(6,0)(0,2){4}{\pspolygon*[linecolor=scacchiera](0,0)(2,0)(2,1)(0,1)(0,0)}
	\multips(2,1)(0,2){4}{\pspolygon*[linecolor=scacchiera](0,0)(4,0)(4,1)(0,1)(0,0)}
	
	\multido{\n=0+1}{9}{
	\psline{-}(\n,0)(\n,8)
	\psline{-}(0,\n)(8,\n)
	\psarc[linewidth=0.6\pslinewidth]{<->}(8,2){0.7}{-70}{70}
	\psarc[linewidth=0.6\pslinewidth]{<->}(6,6){2.7}{-35}{35}
	}
	
	\endpspicture
	\hfill
	\pspicture(0,-0.8)(8,8)
	
	\pspolygon*[linecolor=scacchiera](0,0)(2,0)(2,2)(0,2)(0,0)
	\pspolygon*[linecolor=scacchiera](0,6)(2,6)(2,8)(0,8)(0,6)
	\pspolygon*[linecolor=scacchiera](6,0)(8,0)(8,2)(6,2)(6,0)
	\pspolygon*[linecolor=scacchiera](6,6)(8,6)(8,8)(6,8)(6,6)
	\pspolygon*[linecolor=scacchiera](2,2)(6,2)(6,6)(2,6)(2,2)
	
	\multido{\n=0+1}{9}{
	\psline{-}(\n,0)(\n,8)
	\psline{-}(0,\n)(8,\n)
	}
	
	\endpspicture
	\hfill\mbox{}
\end{center}
}

\renewcommand{\SoluzioniAlternative}{
}

\renewcommand{\Osservazioni}{\TitoloOss{1}
Il primo invariante non bastava per escludere la scacchiera centrale,
nella quale tutte le righe e tutte le colonne contengono 4 caselle
colorate.

Il secondo invariante, applicato alle colonne, bastava invece anche 
per escludere la casella di sinistra, nella quale abbiamo 5 colonne di 
un tipo e 3 colonne di un altro tipo.

\TitoloOss{2} Nel dimostrare che quelli che abbiamo introdotto sono
effettivamente invarianti, abbiamo usato il \emph{principio di
induzione} senza dirlo esplicitamente.  Abbiamo infatti mostrato che
una certa propriet\`{a}, valida nella configurazione iniziale, non
viene alterata dalle mosse possibili, e da questo \`{e} facile dedurre
(formalmente andrebbe dimostrato per induzione sul numero delle mosse
eseguite) che la propriet\`{a} resta valida in \emph{tutte} le
configurazioni successive.

\TitoloOss{3} Ci sono ovviamente molti altri modi, diversi da quello
mostrato, per ottenere la scacchiera di destra.  Non \`{e} comunque
difficile dimostrare che sono necessari almeno 2 scambi di riga ed
almeno 2 scambi di colonna.  }

\renewcommand{\Commenti}{
}